%!TEX root=../../note
\subsection{Implication}
\begin{definition}
    The truth value for $A implies B"$, written symbolically as $A \implies B$, is defined by the truth table 
    \begin{center}
        \begin{tabular}{c|c|c}
        $A$ & $B$ & $A \implies B$ \\
        \hline
        T & T & T \\
        T & F & F \\
        F & T & T \\
        F & F & T
        \end{tabular}
    \end{center}

    We also refer to $A \implies B$ as an implication. 

    Just for fun. Instructor does not talk about it. 
    \begin{equation*}
        A \implies B \equiv \neg (A \wedge \neg B) \equiv \neg A \vee B. 
    \end{equation*}
\end{definition}

Hypothesis: Assume to be true: A

Conclusion: B

\begin{example}
    \begin{equation*}
        \paren{\forall x \in \R}\bracks{\paren{x > 2 } \implies \paren{x^2 > 4}}
    \end{equation*}
\end{example}

\begin{remark}
    Interesting. 
    \begin{equation*}
        (A \vee B) \implies C \equiv (A \implies C) \wedge (B \implies C)
    \end{equation*}
\end{remark}

\subsection{Converse and Contrapositive}
\subsubsection{Converse}
The $B \implies A$ is called the converse of $A \implies B$. 
\begin{center}
    \begin{tabular}{c|c|c|c}
        $A$ & $B$ & $A \implies B$ & $B \implies A$ \\
        \hline
        T & T & T & T \\
        T & F & F & T \\
        F & T & T & F \\
        F & F & T & T
    \end{tabular}
\end{center}
\begin{remark}
    $(A \implies B) \not\equiv (B \implies A)$ (not logical equivalent) 
\end{remark}

\subsubsection{Contrapositive}
The implication $(\neg B) \implies (\neg A)$ is called the contrapositive of $A \implies B$. 
\begin{center}
\begin{tabular}{c|c|c|c|c|c}
$A$ & $B$ & $\neg A$ & $\neg B$ & $A \implies B$ & $(\neg B) \implies (\neg A)$ \\
\hline
T & T & F & F & T & T \\
T & F & F & T & F & F \\
F & T & T & F & T & T \\
F & F & T & T & T & T
\end{tabular}
\end{center}

\begin{remark}
    $(A \implies B) \equiv ((\neg B) \implies (\neg A))$
\end{remark}

An implication is also logically equvalent to another type of compound statement, as shown in the following truth table: 
\begin{center}
\begin{tabular}{c|c|c|c|c}
$A$ & $B$ & $\neg A$ & $A \implies B$ & $\neg A \vee B$ \\
\hline
T & T & F & T & T \\
T & F & F & F & F \\
F & T & T & T & T \\
F & F & T & T & T
\end{tabular}
\end{center}

\begin{remark}
    The negation of an implication is not an implication. 
\end{remark}

\subsubsection{If and Only If}
\begin{definition}
    The truth value for "A if and only if B", written symbolically as $A \iff B$, is defined by the truth table. 

    \begin{center}
    \begin{tabular}{c|c|c}
    $A$ & $B$ & $A \iff B$ \\
    \hline
    T & T & T \\
    T & F & F \\
    F & T & F \\
    F & F & T
    \end{tabular}
    \end{center}

\end{definition}